% !PRINTED 11
\weberpage{11}

%<<CONT
results in the three cases:

\[
\begin{array}{ll}
1) & \dfrac{dx}{\sqrt{-x(1-x)(1-\lambda^2x)}}, \\[10pt]
2) & \dfrac{\lambda\,dx}{\sqrt{x(1-x)(1-\lambda^2x)}}, \\[10pt]
3) & \dfrac{dx}{\sqrt{x(1-x)(1-\lambda^2x)}}.
\end{array}
\]

\begin{center}
\textbf{\S\,4. The \gesperrt{Legendre} Normal Form.}
\end{center}

If the function $f(x)$, which stands under the root sign in the elliptic differential, has real coefficients, then three cases are possible:

1. $f(x)$ has four real roots $x_1$, $x_2$, $x_3$, $x_4$;

2. $f(x)$ has two real roots $x_1$, $x_2$ and a pair of conjugate imaginary roots $x_3$, $x_4$;

3. $f(x)$ has two pairs of conjugate imaginary roots $x_1$, $x_2$ and $x_3$, $x_4$.

In all three cases the elliptic differential can be brought by a real linear transformation into the form:

\weq{1}{\dfrac{dz}{\sqrt{(z^2-\alpha)(z^2-\beta)}},}

in which $\alpha$ and $\beta$ are real (positive or negative) constants.

We determine the linear dependence between the variables $x$ and $z$ in such a way that the following values correspond to one another:

\weq{2}{\begin{array}{ccccc}
x, & x_1, & x_2, & x_3, & x_4, \\
z, & \sqrt{\alpha}, & -\sqrt{\alpha}, & \sqrt{\beta}, & -\sqrt{\beta},
\end{array}}

and obtain a substitution of the form:

\weq{3}{h\,\dfrac{x-x_1}{x-x_2} = \dfrac{z-\sqrt{\alpha}}{z+\sqrt{\alpha}}.}

In order to determine $h$ we set $x = x_3$, $x = x_4$ and correspondingly $z = \sqrt{\beta}$, $z = -\sqrt{\beta}$, and obtain

\[
h\,\dfrac{(x_3x_1)}{(x_3x_2)} = \dfrac{\sqrt{\beta}-\sqrt{\alpha}}{\sqrt{\beta}+\sqrt{\alpha}}, \quad h\,\dfrac{(x_4x_1)}{(x_4x_2)} = \dfrac{\sqrt{\beta}+\sqrt{\alpha}}{\sqrt{\beta}-\sqrt{\alpha}}.
\]

From this, by multiplication and division,

\weq{4}{h = \sqrt{\dfrac{(x_3x_2)(x_4x_2)}{(x_3x_1)(x_4x_1)}}, \quad \dfrac{\sqrt{\alpha}-\sqrt{\beta}}{\sqrt{\alpha}+\sqrt{\beta}} = \sqrt{\dfrac{(x_3x_1)(x_4x_2)}{(x_3x_2)(x_4x_1)}}.}
